\chapter{Despliegue}
\label{chap:despliegue}

\section{Alcance del despliegue documentado}
\label{sec:despliegue-alcance}

El proyecto esta preparado tecnicamente para ejecutarse como aplicacion Django con WSGI o ASGI, variables de entorno, archivos estaticos recolectables, media local y base PostgreSQL fuera de desarrollo. Sin embargo, no hay evidencia de proveedor de hosting, dominio, certificado HTTPS, reverse proxy, servidor de aplicaciones, almacenamiento externo de media ni pipeline CI/CD.

Por tanto, este capitulo describe el proceso recomendado con base en archivos existentes. Todo dato de infraestructura final se marca como PENDIENTE DE DEFINICION PARA PRODUCCION.

\section{Entornos}
\label{sec:despliegue-entornos}

\begin{longtable}{p{3.2cm}p{4.7cm}p{4.8cm}}
\caption{Diferencias entre desarrollo y produccion.}
\label{tab:despliegue-entornos}\\
\toprule
Aspecto & Desarrollo & Produccion \\
\midrule
\variable{DEBUG} & \variable{True} permitido. & \variable{False} obligatorio. \\
Base de datos & SQLite permitido si \variable{DJANGO_DATABASE=sqlite}. & PostgreSQL requerido por configuracion. \\
\variable{SECRET_KEY} & Puede usar clave local de desarrollo. & Debe ser secreto real generado fuera del repositorio. \\
\variable{ALLOWED_HOSTS} & Puede incluir localhost y testserver. & Debe contener dominio o hosts reales. \\
Archivos estaticos & Servidos por \comando{runserver}. & Requieren \comando{collectstatic} y servicio de estaticos. \\
Media & Local en \archivo{media/}. & Estrategia final PENDIENTE DE DEFINICION PARA PRODUCCION. \\
Correo & Consola o SMTP de prueba. & SMTP real validado y confirmado. \\
PDF invitaciones & LibreOffice opcional si no se prueban PDFs. & LibreOffice requerido para conversion DOCX a PDF. \\
\bottomrule
\end{longtable}

\section{Proceso recomendado}
\label{sec:despliegue-proceso}

El flujo recomendado de preparacion es:

\begin{enumerate}
  \item crear entorno virtual de Python;
  \item instalar dependencias desde \archivo{requirements.txt};
  \item crear \archivo{.env} real fuera de control de versiones;
  \item configurar \variable{DJANGO_DEBUG=False};
  \item configurar \variable{DJANGO_SECRET_KEY}, \variable{DJANGO_ALLOWED_HOSTS} y origenes CSRF;
  \item configurar PostgreSQL con variables \variable{POSTGRES_*};
  \item aplicar migraciones;
  \item crear superusuario;
  \item recolectar estaticos;
  \item validar \comando{manage.py check --deploy};
  \item configurar servidor WSGI o ASGI final;
  \item validar SMTP y LibreOffice si se usaran comunicaciones.
\end{enumerate}

\begin{lstlisting}[language=bash,caption={Comandos base de despliegue Django.},label={lst:despliegue-comandos}]
python -m venv .venv
pip install -r requirements.txt
python manage.py migrate
python manage.py createsuperuser
python manage.py collectstatic
python manage.py check --deploy
\end{lstlisting}

\begin{advertencia}
Los comandos de migracion, creacion de superusuario y recoleccion de estaticos modifican el entorno donde se ejecutan. Deben correr en una instancia preparada para despliegue, no durante revision documental ni sobre datos productivos sin respaldo.
\end{advertencia}

\section{Variables obligatorias}
\label{sec:despliegue-variables}

\begin{longtable}{p{4.2cm}p{5.5cm}p{3.2cm}}
\caption{Variables de entorno criticas para despliegue.}
\label{tab:despliegue-variables}\\
\toprule
Variable & Uso & Estado \\
\midrule
\variable{DJANGO_SECRET_KEY} & Clave criptografica de Django. & Obligatoria en produccion. \\
\variable{DJANGO_DEBUG} & Controla modo debug. & Debe ser \variable{False}. \\
\variable{DJANGO_ALLOWED_HOSTS} & Hosts aceptados por Django. & PENDIENTE DE DEFINICION PARA PRODUCCION. \\
\variable{DJANGO_CSRF_TRUSTED_ORIGINS} & Origenes confiables para CSRF. & PENDIENTE DE DEFINICION PARA PRODUCCION. \\
\variable{POSTGRES_DB} & Nombre de base PostgreSQL. & Obligatoria fuera de SQLite. \\
\variable{POSTGRES_USER} & Usuario PostgreSQL. & Obligatoria fuera de SQLite. \\
\variable{POSTGRES_PASSWORD} & Password PostgreSQL. & Obligatoria fuera de SQLite. \\
\variable{POSTGRES_HOST} & Host PostgreSQL. & Obligatoria fuera de SQLite. \\
\variable{POSTGRES_PORT} & Puerto PostgreSQL. & Predeterminado \variable{5432}. \\
\variable{EMAIL_*} & SMTP para comunicaciones reales. & Requerido si hay envios. \\
\variable{LIBREOFFICE_BINARY} & Conversor DOCX a PDF si no esta en PATH. & Requerido para invitaciones PDF. \\
\bottomrule
\end{longtable}

\section{Archivos estaticos, media y documentos}
\label{sec:despliegue-static-media}

La configuracion define \variable{STATIC_ROOT} como \archivo{staticfiles/}, \variable{STATICFILES_DIRS} como \archivo{static/}, \variable{MEDIA_ROOT} como \archivo{media/} y \variable{MEDIA_URL} como \ruta{media/}. En desarrollo, \comando{runserver} puede resolver estaticos. En produccion, el servicio de archivos estaticos queda PENDIENTE DE DEFINICION PARA PRODUCCION.

El almacenamiento de media tambien queda pendiente. La aplicacion puede generar o almacenar archivos asociados a comunicaciones y cargas administrativas; antes de operar con datos reales se debe definir respaldo y retencion para \archivo{media/}.

\section{Comunicaciones en produccion}
\label{sec:despliegue-comunicaciones}

Para envio real se requiere SMTP completo y validado. El proyecto bloquea envios oficiales sin confirmacion. Antes de habilitar envios reales debe probarse:

\begin{itemize}
  \item configuracion SMTP con valores reales;
  \item envio en modo test hacia \variable{test_recipient};
  \item conversion DOCX a PDF con LibreOffice;
  \item ausencia de placeholders en credenciales;
  \item plantillas y destinatarios con datos ficticios o autorizados.
\end{itemize}

\section{Elementos no definidos}
\label{sec:despliegue-pendientes}

\begin{longtable}{p{4.5cm}p{7.8cm}}
\caption{Elementos pendientes para produccion.}
\label{tab:despliegue-pendientes}\\
\toprule
Elemento & Estado \\
\midrule
Proveedor de hosting & PENDIENTE DE DEFINICION PARA PRODUCCION. \\
Dominio final & PENDIENTE DE DEFINICION PARA PRODUCCION. \\
HTTPS y certificados & PENDIENTE DE DEFINICION PARA PRODUCCION. \\
Reverse proxy & PENDIENTE DE DEFINICION PARA PRODUCCION. \\
Servidor WSGI o ASGI final & PENDIENTE DE DEFINICION PARA PRODUCCION. \\
Servicio de static files & PENDIENTE DE DEFINICION PARA PRODUCCION. \\
Servicio de media files & PENDIENTE DE DEFINICION PARA PRODUCCION. \\
Backups de base y media & PENDIENTE DE DEFINICION PARA PRODUCCION. \\
Monitoreo de errores & PENDIENTE DE DEFINICION PARA PRODUCCION. \\
CI/CD & PENDIENTE DE DEFINICION PARA PRODUCCION. \\
\bottomrule
\end{longtable}

\section{Recomendaciones para futuras versiones}
\label{sec:despliegue-recomendaciones}

\begin{itemize}
  \item Crear una guia de produccion cuando se confirme proveedor, dominio, HTTPS, servidor de aplicacion y estrategia de static/media.
  \item Agregar checklist de despliegue versionado con responsables, fecha y resultado de pruebas.
  \item Validar \comando{manage.py check --deploy} en el entorno final antes de abrir el sistema al publico.
  \item Documentar procedimiento de rollback de version y restauracion de base de datos.
  \item Automatizar despliegue solo despues de definir repositorio remoto, secretos, migraciones y backups.
\end{itemize}
